\section{Introduction}

Paper~10~\cite{paper10adaptive} evaluated a one-threshold magnitude
change-point detector for online recalibration of HygienePrio's
scorer weights at the single pre-registered $\tau = 0.05$. The
outcome was mixed: H1 supported (no hazard at $K = 200$), H2 rejected
($5.3$~pp penalty at $K = 50$), H3 rejected (non-monotone firing in
capacity), H4 supported (marginal beat over static gate). Paper~10's
discussion argued the H1 success and H2 cost are parameterised by
$\tau$ and form a Pareto curve over the threshold axis.

This paper tests that argument. We sweep $\tau \in \{0.02, 0.03,
0.05, 0.075, 0.10\}$ on the same $K \in \{50, 100, 200\}$,
$\lambda = 3$ grid as Paper~10 and ask four pre-registered questions:
does the K=200 hazard re-emerge at small $\tau$ (H1); is the K=50
cost a monotone function of $\tau$ (H2); does any $\tau$ satisfy
both Paper~10 tolerances simultaneously (H3); is the firing fraction
monotone in $\tau$ (H4)?

\subsection{Contributions}

\begin{itemize}
\item \textbf{C1 --- A pre-registered $\tau$ sweep.} The first
pre-registered sensitivity analysis of the Paper~10 adaptive detector
across an order-of-magnitude $\tau$ grid; 11{,}250 frozen result rows
plus a 90-row change-point decision log.

\item \textbf{C2 --- H3 rejected: simple detector cannot reach the
feasible region.} No $\tau$ in the swept grid simultaneously
satisfies both Paper~10 tolerances. The K=50 cost is between $-4.1$
and $-5.3$~pp at every $\tau$, always exceeding the $-2.0$~pp
ceiling. The simple magnitude detector is fundamentally insufficient;
richer detectors are necessary.

\item \textbf{C3 --- H1 is robust across the $\tau$ grid.} The
worst per-window $K = 200$ adaptive--fixed delta is exactly $0.0$ at
every $\tau$ tested. The hazard-avoidance property of the
detector is not specific to Paper~10's $\tau = 0.05$ choice but
holds across an order-of-magnitude variation in the threshold ---
a stronger H1 statement than Paper~10 reported.

\item \textbf{C4 --- H2 and H4 protocol-sign reversal, honestly
reported.} The pre-registration predicted firing fraction would be
monotone non-decreasing in $\tau$. The data reveals the opposite:
$\rho(\tau, \text{fire frac}) = -0.97$ ($p = 0.005$). Larger $\tau$
requires a larger shift to trigger and therefore fires less; the
protocol mathematically inverted the direction. H2 inherited the
sign error. We report the rejections honestly per the protocol's
stop rules and use the data direction as the discussion's framing.

\item \textbf{C5 --- A negative operating-point result.} The
Paper~10 question ``does an operating point exist for the simple
detector that satisfies both tolerances?'' has a negative answer in
this grid. The answer requires either a richer detector family
(CUSUM, Bayesian) or relaxing one of the two Paper~10 tolerances.
\end{itemize}

\subsection{Relationship to Prior Papers}

This paper is the eleventh in the VulnPrio research sequence. It
closes the simple-detector branch opened by Paper~10 with a sharp
negative-feasibility result and opens the richer-detector branch as
the necessary next step.
